Un alternateur simple est formé d'une bobine rectangulaire de $N = \v{N}$ spires, de dimensions \qty{\v{a}}{cm} sur \qty{\v{b}}{cm}, qui tourne à \qty{\v{f}}{Hz} dans un champ magnétique uniforme de \qty{\v{B}}{T}. L'axe de rotation est perpendiculaire au champ. À $t = 0$, le plan de la bobine est perpendiculaire au champ.

\begin{enumerate}[label=\alph*)]
    \item Déterminez la vitesse angulaire de la bobine.
    \item Déterminez la force électromotrice maximale produite.
    \item Combien de temps après $t = 0$ la force électromotrice atteint-elle pour la première fois la moitié de sa valeur maximale ?
\end{enumerate}
